\documentclass[10pt,aspectratio=169]{beamer}

% -------------------------------------------------
% Theme (Solarized Light)
% -------------------------------------------------
\usetheme{metropolis}
\metroset{
  progressbar=none,
  numbering=fraction
}
\setbeamertemplate{navigation symbols}{}

% -------------------------------------------------
% Fonts -- requires LuaLaTeX or XeLaTeX
% -------------------------------------------------
\usepackage{fontspec}
\setmainfont{Latin Modern Roman}
\setsansfont{Latin Modern Sans}
\setmonofont{Latin Modern Mono}

% -------------------------------------------------
% Colors: Solarized Light
%   Base3  #fdf6e3  background
%   Base2  #eee8d5  background highlight
%   Base1  #93a1a1  secondary content
%   Base00 #657b83  body text
%   Base01 #586e75  emphasis / headings
%   Blue   #268bd2  primary accent
%   Cyan   #2aa198  secondary accent
% -------------------------------------------------
\usepackage{xcolor}
\definecolor{SolBase3} {HTML}{fdf6e3}
\definecolor{SolBase2} {HTML}{eee8d5}
\definecolor{SolBase1} {HTML}{93a1a1}
\definecolor{SolBase00}{HTML}{657b83}
\definecolor{SolBase01}{HTML}{586e75}
\definecolor{SolBlue}  {HTML}{268bd2}
\definecolor{SolCyan}  {HTML}{2aa198}

\setbeamercolor{background canvas}{bg=SolBase3}
\setbeamercolor{normal text}{fg=SolBase00}
\setbeamercolor{alerted text}{fg=SolBlue}

\setbeamercolor{frametitle}{bg=SolBase3, fg=SolBase01}
\setbeamercolor{title}{fg=SolBase01}
\setbeamercolor{subtitle}{fg=SolBase00}

\setbeamercolor{progress bar}{fg=SolBlue}
\setbeamercolor{block title}{fg=SolBase3, bg=SolBase01}
\setbeamercolor{block body}{bg=SolBase2, fg=SolBase00}

% Optional refinements
\setbeamercolor{section title}{fg=SolBase01}
\setbeamercolor{footline}{fg=SolBase1}
\setbeamercolor{item}{fg=SolBlue}
\setbeamercolor{subitem}{fg=SolCyan}

% -------------------------------------------------
% Packages
% -------------------------------------------------
\usepackage{graphicx}
\usepackage{booktabs}
\usepackage{amsmath}
\usepackage{hyperref}
\hypersetup{
  colorlinks=true,
  linkcolor=SolBlue,
  urlcolor=SolCyan,
  citecolor=SolBlue
}

% -------------------------------------------------
% Typography and spacing tweaks
% -------------------------------------------------
\usefonttheme{professionalfonts}
\setlength{\parskip}{0.4em}
\linespread{1.05}

% Slightly more breathing room in tables
\renewcommand{\arraystretch}{1.15}

% Cleaner captions if figures/tables are used
\setbeamerfont{caption}{size=\small}
\setbeamercolor{caption name}{fg=SolBase01}

% -------------------------------------------------
% Useful macros
% -------------------------------------------------
\newcommand{\highlight}[1]{\textcolor{SolBlue}{\textbf{#1}}}
\newcommand{\source}[1]{\vspace{0.5em}\hfill{\scriptsize\textcolor{SolBase1}{#1}}}

% -------------------------------------------------
% Title information
% -------------------------------------------------
\title{Title of the Talk}
\subtitle{A concise and informative subtitle}
\author{Your Name}
\institute{Your Institution}
\date{\today}

\begin{document}

% -------------------------------------------------
% Title slide
% -------------------------------------------------
\begin{frame}
  \titlepage
\end{frame}

% -------------------------------------------------
% Example section
% -------------------------------------------------
\section{Motivation}

\begin{frame}{Why this question matters}
\begin{itemize}
  \item Regions do not diversify randomly.
  \item Existing capabilities shape which new activities can emerge.
  \item The interesting question is not whether change happens, but \highlight{how}.
\end{itemize}
\end{frame}

\begin{frame}{A visual slide should look like this}
  \centering
  \includegraphics[width=0.82\textwidth]{figures/example.pdf}
  
  \source{Source: Author's calculations}
\end{frame}

\section{Main result}

\begin{frame}{Related capabilities predict regional entry}
\begin{block}{Main takeaway}
Regions are more likely to enter industries that are cognitively and organizationally related to their existing activities.
\end{block}
\end{frame}

\begin{frame}{Illustrative table}
\centering
\begin{tabular}{lcc}
\toprule
Variable & Model 1 & Model 2 \\
\midrule
Relatedness & 0.245*** & 0.198*** \\
Controls & Yes & Yes \\
Fixed effects & No & Yes \\
\bottomrule
\end{tabular}

\vspace{0.5em}
{\small Standard errors omitted here because this is a slide, not a punishment.}
\end{frame}

\section{Conclusion}

\begin{frame}{Conclusion}
\begin{itemize}
  \item Diversification is structured by existing capabilities.
  \item Mechanisms matter more than slogans.
  \item Policy works better when it respects regional capability constraints.
\end{itemize}
\end{frame}

\end{document}